\documentclass[12pt, a4paper]{article}

% --- Packages ---
\usepackage[utf8]{inputenc}
\usepackage[T1]{fontenc}
\usepackage{lmodern}
\usepackage[margin=1in]{geometry}
\usepackage{setspace}
\usepackage{amsmath, amssymb}
\usepackage{booktabs}
\usepackage{tabularx}
\usepackage{hyperref}
\usepackage{natbib}
\usepackage{graphicx}
\usepackage{float}
\usepackage{url}
\usepackage{xcolor}

% --- Hyperref setup ---
\hypersetup{
  colorlinks=true,
  linkcolor=blue!60!black,
  citecolor=blue!60!black,
  urlcolor=blue!60!black
}

% --- Spacing ---
\onehalfspacing

% --- Title ---
\title{\textbf{From Labor to Attention: A Token-Economic Framework for the AI Era}\thanks{Available at SSRN: \url{https://papers.ssrn.com/sol3/papers.cfm?abstract_id=6486722}}}
\author{Peiya Sun\thanks{Independent Researcher, cPen Network. Contact: peter.sun@cpen.io.}}
\date{March 2026}

\begin{document}

\maketitle

\begin{abstract}
\noindent
Modern economies distribute income and legitimacy primarily through labor. As AI and automation expand across cognitive, physical, and administrative work, that arrangement becomes more fragile \citep{frey2017, acemoglu2025, eloundou2024}. The same forces that weaken labor-centered income also increase the economic centrality of human attention: AI reduces the demand for human labor while increasing the supply of goods and content competing for finite human attention \citep{simon1971, goldhaber1997}. These two dynamics have been studied separately but not connected.

This paper bridges that gap with three contributions. First, the Attention Triad distinguishes attention contribution, aggregation, and monetization as separate economic functions, making visible the contribution gap in current attention markets. Second, the attention-centered economy is proposed as a structural parallel to the labor-centered economy, arguing that the forces weakening labor-based inclusion simultaneously strengthen the case for attention-based compensation. Third, the Attention-Token Economic Model (ATEM) specifies a governance-conscious institutional framework for compensating verified human attention, illustrated through a tokenized architecture with revenue-backed sinks. A formal model identifies viability conditions including critical mass thresholds, and a sensitivity analysis demonstrates that viability is dominated by network effects. These mechanisms are presented as institutional design tools whose legitimacy depends on demanding conditions of privacy, proof-of-humanity, governance, and restraint.

\medskip
\noindent\textbf{Keywords:} attention economy, platform economics, digital labor, token economics, automation, artificial intelligence, value distribution, governance

\medskip
\noindent\textbf{JEL Classification:} O33, D83, L86, P16, J01
\end{abstract}

\newpage
\tableofcontents
\newpage

%% ========================================================================
\section{Introduction}
\label{sec:introduction}
%% ========================================================================

Modern economies remain labor-centered: wages, social insurance, professional identity, and access to credit and housing all depend heavily on stable employment. Artificial intelligence and automation place that arrangement under increasing pressure, not only through outright job displacement but through the slower erosion of employment's reliability as a channel of economic participation \citep{brynjolfsson2014, ford2015}. At the same time, platform economies have made human attention one of the most actively monetized inputs in the contemporary economy, yet the value generated by attention markets flows overwhelmingly to platforms, creators, advertisers, and merchants, while ordinary users whose attention makes those markets possible receive little direct compensation \citep{terranova2012, zuboff2019}.

These two observations have been studied in separate literatures. Attention economists analyze how attention creates value in platform-mediated environments but do not connect their analysis to labor fragility. Post-work scholars analyze labor displacement and propose fiscal responses but do not identify attention compensation as a possible complementary mechanism. Yet the underlying dynamics are connected: the same forces (AI, automation, platform mediation) that weaken labor-centered income simultaneously increase the economic centrality of human attention.

This paper bridges that gap and makes three linked contributions:

\begin{enumerate}
  \item \textbf{The Attention Triad} (Section~\ref{sec:triad}): An analytic framework that distinguishes contribution, aggregation, and monetization as separate economic functions in attention markets, making the contribution gap visible as a structural rather than incidental feature.

  \item \textbf{The Attention-Centered Economy} (Section~\ref{sec:ace}): A structural parallel to the labor-centered economy, with a direct component-to-component mapping. The paper argues that the attention-centered economy may emerge as a complementary and gradually expanding layer of economic participation as AI and automation weaken labor-based inclusion.

  \item \textbf{ATEM} (Sections~\ref{sec:atem}--\ref{sec:limits}): A governance-conscious institutional framework for attention compensation, specifying design requirements, a tokenized architecture, a simple formal model of token dynamics, and the structural conditions under which such a system may be viable.
\end{enumerate}

The paper proceeds as follows. Section~\ref{sec:labor} explains why labor-centered income remains institutionally central and why AI makes that arrangement more fragile. Section~\ref{sec:triad} introduces the Attention Triad. Section~\ref{sec:ace} proposes the attention-centered economy as a structural parallel. Section~\ref{sec:atem} defines ATEM, presents its four core components and tokenized architecture, and positions it against alternative responses. Section~\ref{sec:design} states the design requirements. Section~\ref{sec:model} develops a simple formal model, magnitude estimate, and sensitivity analysis of critical mass. Section~\ref{sec:limits} analyzes limits, risks, and governance problems. Section~\ref{sec:conclusion} concludes with policy implications.


%% ========================================================================
\section{Labor-Centered Income and Its Emerging Fragility}
\label{sec:labor}
%% ========================================================================

This paper uses the term \textbf{labor-centered economy} to describe the current distributive arrangement in which labor serves as the primary channel through which ordinary people access income, social protection, and economic legitimacy. The term does not mean that labor is the only source of value, or that capital, institutions, and technology are secondary. It means that the ordinary person's access to income and social legitimacy is still mediated mainly through work. This is why employment shocks cascade beyond wages: they affect healthcare access, retirement security, social standing, household stability, and political trust.

That labor-centered structure becomes fragile when technology changes the relation between production and human work. Earlier waves of mechanization displaced particular tasks while often creating new forms of employment elsewhere. The current wave of AI and automation raises a different concern: it affects not only routine physical tasks but also clerical, analytical, communicative, and creative tasks that long served as refuges for displaced labor. \citet{frey2017} estimated that 47\% of US employment faces high risk of automation. More recent assessments confirm and extend the concern: \citet{eloundou2024} find that over 46\% of jobs could have more than half their tasks affected by large language models when complementary software is considered, while \citet{acemoglu2025} estimates modest aggregate productivity gains but warns that AI is likely to widen the gap between capital and labor income. The World Economic Forum \citeyearpar{wef2025} projects that 92 million jobs will be displaced by 2030, even as 170 million new roles emerge, implying massive structural reallocation. The scope of vulnerable occupations extends well beyond manufacturing into cognitive and service work \citep{brynjolfsson2014, ford2015}.

The displacement threat should be understood in institutional rather than purely numerical terms. If AI systems absorb a growing share of high-volume cognitive work, firms may need fewer entry-level workers, fewer support workers, and fewer mid-skill knowledge workers to produce the same output. Remaining jobs may become more precarious, more tightly surveilled, or more weakly compensated. Even where total employment does not collapse, the social contract organized around labor may weaken if stable and adequately paid work becomes more unevenly distributed.

Much of the policy debate responds from the fiscal side. Proposals such as Universal Basic Income \citep{standing2017, yang2018} address real needs, and empirical evidence shows some positive effects on well-being \citep{kangas2019}. But UBI and similar proposals face challenges in funding, political feasibility, and sustainability, and they typically operate through ex post fiscal redistribution detached from the sites where value is produced. These approaches are valuable but incomplete. They leave open a further question: whether the contemporary economy already contains other human contributions that are scarce, valuable, and systematically undercompensated.

The most obvious candidate is attention.


%% ========================================================================
\section{The Attention Triad}
\label{sec:triad}
%% ========================================================================

\subsection{Attention as a Scarce Economic Input}
\label{sec:triad-scarce}

Attention has long been recognized as a scarce economic resource. \citet{simon1971} argued that information abundance produces attention scarcity. \citet{goldhaber1997} treated attention as a central currency of networked environments. \citet{davenport2001} and \citet{lanham2006} extended the point into business, media, and institutional settings. \citet{wu2016} traced the historical commodification of attention from newspapers through television to digital platforms. \citet{citton2017} generalized the framing further, and \citet{heitmayer2024} proposes a ``second wave'' of attention economics where attention functions as a ``universal symbolic currency.''

What makes that literature newly urgent is that platform economies have operationalized the insight at massive scale. Platforms do not merely host content. They organize competition for human attention. Search engines, social media platforms, video services, app ecosystems, and digital marketplaces all depend on attracting, holding, routing, and converting user attention. In many cases, the primary revenue model is straightforward: users provide attention, platforms aggregate it into addressable audiences, and advertisers or merchants pay for access to those audiences \citep{rochet2003, wu2016}.

The economic evidence supports the growing scarcity value of attention. Customer acquisition costs in e-commerce have risen significantly while the cost of producing and listing products has fallen steadily \citep{simplicitydx2022}. \citet{teixeira2014} measured a sevenfold real increase in the cost of one minute of consumer attention to advertising over two decades. The economics are inverting: producing goods becomes cheaper while securing human attention becomes more expensive.

\subsection{Three Functions of Attention Markets}
\label{sec:triad-functions}

Once attention is treated as an economically central input, three distinct functions become analytically separable.

\begin{table}[H]
\centering
\caption{Three functions of attention markets}
\label{tab:triad}
\begin{tabularx}{\textwidth}{lXX}
\toprule
\textbf{Function} & \textbf{Description} & \textbf{Typical actors} \\
\midrule
Attention Contribution & Supplying scarce human attention through viewing, reading, comparing, browsing, evaluating, or otherwise engaging & Users, audiences, participants \\
\addlinespace
Attention Aggregation & Attracting, organizing, and directing distributed attention into coherent channels, communities, or audiences & Platforms, creators, curators, publishers \\
\addlinespace
Attention Monetization & Converting aggregated attention into revenue through advertising, subscriptions, commerce, or other paid access models & Advertisers, merchants, service providers, platforms \\
\bottomrule
\end{tabularx}
\end{table}

This paper refers to these three functions collectively as the \textbf{Attention Triad}. The Triad does not require that these roles be occupied by separate people or firms. Platforms often aggregate and monetize at the same time. A creator aggregates attention and may monetize it through sponsorships. An ordinary user primarily contributes attention but could also create content. What the Triad requires is that these functions be analytically distinguishable, even when a single entity performs more than one.

\subsection{Contribution Is Not the Same as Creation}
\label{sec:triad-distinction}

Treating attention contribution as economically meaningful does not mean treating every contribution as equivalent. Creating content, building infrastructure, investing capital, maintaining systems, and converting engagement into revenue involve distinct forms of effort, risk, and skill. The Triad does not erase those differences. It argues that contributors remain participants in value creation even when they are not authors, entrepreneurs, or owners.

That distinction matters because contemporary attention markets usually compensate only successful aggregators or monetizers. The contributor appears only as a user, audience, or data source. But the scarcity of attention means the contributor cannot be reduced to background context. Aggregators and monetizers depend on contributors in the same way that audience markets depend on actual viewers and shoppers rather than simulated ones.

\subsection{The Contribution Gap}
\label{sec:triad-gap}

Critical media and platform scholarship has shown that the attention economy is distributively asymmetric. \citet{smythe1977} described audiences as the commodity sold to advertisers. \citet{terranova2000, terranova2012} treated digital participation as free labor. \citet{fuchs2014} connected user participation to broader patterns of exploitation, while \citet{zuboff2019} showed how behavioral extraction became a source of commercial prediction and control. \citet{srnicek2017} analyzed how platforms systematically capture value from user interactions while externalizing costs.

These traditions differ in emphasis, but they agree on a central point: users are not economically passive simply because they are not employees. This paper calls the resulting asymmetry the \textbf{contribution gap}: contributors supply the scarce input on which attention markets depend, yet they are usually outside the direct distributive loop.

A natural objection is that contributors are already compensated through free services: search, email, social media, and cloud storage. On this view, the market-clearing price of user attention is exactly the value of the services received, and there is no gap to close. This paper offers two responses. First, the ``free services'' framing treats the exchange as bilateral (attention for services), whereas the Triad reveals a trilateral structure: contributors supply attention, aggregators organize it, and monetizers pay for access. The revenue from monetizers substantially exceeds the cost of providing free services to contributors (major platforms routinely operate at 60--80\% gross margins), and the surplus accrues to aggregators and monetizers rather than to the contributors whose scarcity makes the market possible. Second, the ``free services'' exchange is opaque, non-negotiable, and non-contestable: users cannot observe what their attention generates, bargain for a different share, or exit without losing access to essential digital infrastructure. The contribution gap is not a claim that contributors receive nothing. It is a claim that the current distributive arrangement lacks the transparency and contestability that would characterize a fair exchange.

\subsection{Relation to the Data-as-Labor Literature}
\label{sec:triad-dal}

The argument that user contributions deserve compensation has a parallel in the ``data as labor'' literature. \citet{posner2018} argue that data producers should be treated as workers and compensated for the data they generate. \citet{arrieta2018} formalize this as a policy question: if data is labor, users should receive wages for producing it.

These proposals and ATEM address different dimensions of the Attention Triad. In the Triad's terms, data generation is primarily a byproduct of the aggregation layer: platforms collect behavioral data as they organize and direct user attention. The data-as-labor literature is therefore arguing that aggregators should compensate users for what they extract during the aggregation process. ATEM, by contrast, addresses the contribution layer: it argues that users should be compensated for supplying the scarce input (attention) upstream of aggregation, regardless of whether data is also collected.

This distinction matters for institutional design. Data-as-labor proposals require measuring and pricing individual data contributions, which faces formidable valuation challenges: how much is one user's browsing data worth in isolation versus as part of a training set? Attention compensation operates through simpler proxies (presence, engagement, task completion) tied to the act of contributing the scarce resource itself, without requiring fine-grained data accounting.

The two approaches are complementary rather than competing, precisely because they target different functions in the value chain. A platform could, in principle, compensate users both for the data extracted during aggregation and for the attention supplied as contribution. But the institutional mechanisms would differ, because the underlying economic properties differ: data is non-rivalrous and accumulable once collected, while attention is rivalrous and finite in the moment of supply. ATEM focuses on the latter because it is the scarcity of attention, not of data, that makes the contribution gap structurally significant.


%% ========================================================================
\section{The Attention-Centered Economy}
\label{sec:ace}
%% ========================================================================

\subsection{A Structural Parallel}
\label{sec:ace-parallel}

The labor-centered economy is not simply an economy that uses labor. It is a distributive arrangement in which labor serves as the primary channel through which ordinary people access income, social protection, and economic legitimacy. Workers supply the scarce input (labor), firms organize it into productive output, wages compensate it, labor markets match supply and demand, and labor law protects participants.

This paper proposes that an analogous structure is emerging around attention. The structural parallel can be stated precisely through a direct component-to-component mapping:

\begin{table}[H]
\centering
\caption{Structural parallel between labor-centered and attention-centered economies}
\label{tab:parallel}
\begin{tabularx}{\textwidth}{lXX}
\toprule
\textbf{Dimension} & \textbf{Labor-Centered Economy} & \textbf{Attention-Centered Economy} \\
\midrule
Scarce input & Human labor & Human attention \\
Input suppliers & Workers & Contributors (users, audiences) \\
Organizers & Firms, employers & Platforms, aggregators \\
Value conversion & Merchants, distributors & Monetizers (advertisers, merchants) \\
Compensation & Wages, benefits & Token rewards, attention-based claims \\
Output measure & Productivity & Engagement value, conversion \\
Market mechanism & Labor markets & Attention markets \\
Protective framework & Labor law, unions, social insurance & Governance, proof-of-humanity, contestability \\
Distributive channel & Employment & Verified attention contribution \\
\bottomrule
\end{tabularx}
\end{table}

Each row maps a functional component of the labor-centered economy to its structural equivalent in the attention-centered economy. The parallel is not perfect; attention markets are younger and less institutionally developed, and the protective frameworks for attention contributors remain rudimentary or nonexistent. But the structural correspondence is precise enough to ground a serious claim: the attention-centered economy is not a speculative projection but a description of functional relationships that already exist in platform-mediated markets, with the critical difference that the contribution dimension remains uncompensated.

The Attention Triad maps directly onto the supply, organization, and conversion rows of this table. Contributors are the workers of the attention economy. Aggregators are its organizers. Monetizers are its merchants. The contribution gap is the structural equivalent of an industrial economy in which workers supply the scarce input but receive no wages.

\subsection{The Bridging Argument: Why the Same Forces Create Both Problems}
\label{sec:ace-bridging}

The most important feature of this parallel is that the same forces drive both dynamics simultaneously. AI and automation weaken the labor-centered economy by reducing the demand for human labor in production. At the same time, they strengthen the attention-centered economy by increasing the supply of goods, services, and content that compete for finite human attention. As production costs fall across sectors, the bottleneck shifts from what can be produced to what humans can notice, evaluate, and choose among.

This connection is not incidental. It follows from the same technological trajectory. AI systems that automate cognitive work (reducing demand for human labor) also generate more content, products, and services (increasing demand for human attention to discover and evaluate them). The labor-centered economy weakens and the attention-centered economy strengthens as consequences of the same underlying force.

This is the bridging insight that connects two literatures that have developed largely in isolation. Attention economists \citep{simon1971, goldhaber1997, wu2016} analyze how attention creates value but do not connect their analysis to the institutional weakening of labor-centered income. Post-work scholars \citep{standing2017, brynjolfsson2014, ford2015} analyze labor displacement but do not identify attention compensation as a possible complementary mechanism. Yet the conditions that make labor less secure are the same conditions that make attention more valuable.

\subsection{Scope and Limits of the Parallel}
\label{sec:ace-limits}

Several clarifications are necessary. For the foreseeable future, the attention-centered economy is a complement to the labor-centered economy, not a replacement. Labor remains central to production and continues to anchor most people's income. The near-term claim is that in specific sectors where attention is already directly monetized, an attention-centered distributive layer may emerge alongside labor-based arrangements.

The attention-centered economy is not a claim that all economic activity becomes attention-based. Manufacturing, agriculture, infrastructure, healthcare, and many other domains remain organized primarily around labor, capital, and material inputs. The attention-centered economy applies most directly to platform-mediated, advertising-supported, and digitally intermediated domains, though the scope of those domains continues to expand.

In the longer run, the relationship between the two economies depends on the pace and scope of automation. If AI and robotics eventually automate the large majority of human labor, the attention-centered economy could shift from a complementary layer to a more primary one. This paper does not predict that outcome, but the structural parallel is designed to hold under both scenarios.


%% ========================================================================
\section{ATEM: Definition and Architecture}
\label{sec:atem}
%% ========================================================================

\subsection{Definition}
\label{sec:atem-def}

The Attention-Token Economic Model (ATEM) is a conceptual and institutional framework for compensating verified human attention contribution in domains where attention is monetizable. It is this paper's proposed response to the contribution gap identified by the Attention Triad (Section~\ref{sec:triad}) and the distributive opportunity identified by the attention-centered economy framework (Section~\ref{sec:ace}).

The Attention Triad is the analytic foundation of ATEM, not a synonym for it: the Triad explains the functional structure of attention markets, while ATEM specifies a possible distributive response to that structure. A tokenized layer is one possible implementation of ATEM rather than the definition of ATEM itself. Different implementations may vary in their specific choices while adhering to the fundamental principles described here.

\subsection{The Four Core Components}
\label{sec:atem-components}

ATEM operates through four core components that map directly onto the Attention Triad. Each component addresses one functional dimension of the attention economy, and together they form a cycle in which value flows from attention contribution through monetization and back to contributors.

\textbf{Attention Capture.} The first component recognizes and validates human attention contribution. Contributors (the supply side of the Triad) direct their scarce attention toward content, products, or information, and the system registers that contribution through Attention Capture Mechanisms (ACMs). ACMs span a spectrum of effort and complexity. At the minimal end, a simple check-in or presence confirmation takes seconds and requires no special skill, ensuring universal accessibility. Further along the spectrum, natural engagement behaviors such as browsing, evaluating products, posting, or commenting capture attention through ordinary platform use. At the higher end, task-based mechanisms such as surveys, content evaluation, or data labeling capture deeper attention and generate higher-value signals. The defining requirement is that any ACM captures genuine human attention, verified through proof-of-humanity validation. Any implementation must include at least one minimal-effort ACM to ensure that no one is excluded for lack of content, skills, or sustained engagement time.

\textbf{Attention Monetization.} The second component converts aggregated attention into revenue. This corresponds to the monetization dimension of the Triad: advertisers, merchants, and service providers pay for access to verified human attention. Revenue sources may include advertising, subscription fees, marketplace commissions, transaction fees, or other monetization channels. The key principle is that revenue is generated from real economic demand for human attention, not from speculative token trading.

\textbf{Token Distribution.} The third component distributes tokens to contributors in proportion to their verified attention contribution. This is the mechanism that closes the contribution gap: contributors receive a direct claim on the value their attention helps generate, rather than only indirect benefits through free services. Distribution must balance sufficient rewards to incentivize participation, controlled issuance to prevent inflation, and fair allocation across different forms and intensities of attention contribution.

\textbf{Token Absorption.} The fourth component removes tokens from circulation using system revenue. A portion of the revenue generated through attention monetization is used to buy back tokens from the open market, absorbing them from circulating supply. This creates a direct link between contributor rewards and real economic demand: the tokens participants receive are backed by actual revenue, not by detached issuance. Token absorption may also occur through other mechanisms: tokens spent on in-platform services, governance staking, promotional spending, or transaction fees.

Token absorption through revenue-backed buybacks is the strongest mechanism claim in this paper. Many token systems fail precisely by distributing claims faster than they generate durable reasons for others to want them. Revenue-backed token sinks connect the compensation layer to buyer-side demand rather than asking participants to trust pure issuance or speculative appreciation. This mechanism is not sufficient on its own: a system can still fail if buyer demand is weak, if fraud overwhelms validation, if governance is captured, or if volatility makes rewards unusable. Revenue-backed token sinks are one necessary design discipline among several, not a proof of self-stabilization.

Together, these four components form a self-reinforcing cycle: attention capture generates the input, monetization converts it to revenue, distribution compensates contributors, and absorption closes the loop by linking token value to real demand. As long as monetizers continue to pay for access to human attention, the cycle can sustain itself: revenue funds absorption, absorption supports token value, token value incentivizes participation, and participation generates the attention that monetizers pay for.

\subsection{Why a Tokenized Layer}
\label{sec:atem-tokens}

The case for tokens in ATEM is practical rather than ideological. In principle, attention compensation could be delivered through fiat payments, platform credits, cooperative accounting systems, or public transfers. ATEM's normative force does not depend on claiming that tokens are uniquely just; it depends on whether contributors receive fairer, more transparent, and more contestable claims on value.

A tokenized layer is relevant because programmable digital assets can coordinate adoption, incentives, and distribution in ways that conventional payment rails do not handle as cleanly \citep{catalini2018, cong2021}. Specifically: programmable assets can encode distribution and absorption rules directly; they can support very small, frequent claims without payroll-like structures; they can make part of the distributive side visible and auditable; and they can link ecosystem performance to participant rewards through transparent on-chain mechanisms.

None of these advantages is automatic. Tokens also create volatility, liquidity risk, governance conflicts, and regulatory uncertainty. Tokenized infrastructure is therefore best understood as one implementation option that may help under certain conditions, especially when transparency, programmability, and micropayment logic matter.

\subsection{Positioning}
\label{sec:atem-position}

ATEM should be understood as a contingent institutional response within specific sectors, not as a wholesale replacement for labor markets or welfare states. This section clarifies how ATEM relates to two prominent alternatives: fiscal redistribution (including UBI) and existing platform reward systems.

\textbf{Against pure fiscal redistribution.} Universal Basic Income and similar proposals respond to labor fragility through ex post transfers funded by taxation. That is an essential approach, but it operates at a distance from the sites where value is produced. ATEM attempts something different: building a distributive claim inside markets that already monetize human attention. Rather than taxing platform revenue and redistributing it through public transfers, ATEM embeds compensation into the economic relationship between contributors and monetizers. This does not make UBI unnecessary. It means that attention compensation could complement fiscal redistribution by operating within the market where the undercompensation occurs, rather than only correcting for it afterward.

\textbf{Against current platform reward systems.} Existing platforms do compensate users in various ways: creators and sellers receive revenue-sharing, and ordinary users may receive discounts, loyalty points, promotional credits, or sweepstake entries. These mechanisms acknowledge that user attention has value. But they are typically ad hoc, discretionary, and selective rather than systematic. Discounts and promotions are marketing tools designed to drive specific behaviors, not general compensation for attention contribution. They are offered intermittently, to targeted segments, under opaque rules set unilaterally by platforms, and they can be withdrawn at any time without recourse. Some token-based systems have begun rewarding users for specific actions such as viewing advertisements, but these typically compensate one narrow form of attention rather than the broader range of contributions that generate value across a platform's ecosystem. ATEM differs in three respects. First, it treats attention compensation as a general, regular entitlement for all verified contributors, not a reward for a single action. Second, it makes revenue linkage a structural requirement rather than a discretionary feature. Third, it requires governance, proof-of-humanity, and contestability as design foundations rather than afterthoughts.

\textbf{Scope.} ATEM is most plausible in domains where attention is already directly monetized: platform-mediated media, digital marketplaces, advertising-supported ecosystems, social commerce, and evaluation or curation environments. In those settings, the necessary infrastructure for measuring, aggregating, and monetizing attention already exists. What is missing is the distributive link back to contributors.


%% ========================================================================
\section{Design Requirements}
\label{sec:design}
%% ========================================================================

Any attempt to compensate attention contributors must satisfy stricter conditions than most current platform incentive systems do. The challenge is not simply to distribute tokens or points. The challenge is to do so without collapsing into manipulation, surveillance, or speculation.

\textbf{Real-demand linkage.} The system must connect participant rewards to actual monetizable demand. Pure issuance, untethered from real revenue, invites speculative boom-bust dynamics.

\textbf{Accessibility without high exclusion.} Participation should not require capital, professional skill, or advanced technical knowledge. At the same time, the system should not pretend that all forms of attention are equally valuable. A workable design must preserve broad access without collapsing all contributions into an undifferentiated mass.

\textbf{Measurement humility.} Attention cannot be measured in a pure philosophical sense. What systems can observe are proxies: presence, liveness, response, task completion, persistence, and sometimes downstream conversion. A serious framework must be modest about what it measures and honest about the gap between inner attention and operational signals.

\textbf{Proof-of-humanity with privacy restraint.} Any system that rewards participation will face pressure from bots, farms, and increasingly capable AI agents. Some form of proof-of-humanity or anti-sybil protection is necessary. But stronger verification often means more intrusive data collection. If validation becomes too invasive, the framework undermines part of its own normative case. Existing approaches, from Worldcoin's biometric iris scanning to Gitcoin Passport's credential aggregation, illustrate the spectrum of tradeoffs. No current solution fully resolves the tension between robust verification and minimal data collection.

\textbf{Transparency, contestability, and governance.} Participants should be able to understand how rewards are allocated, how tokens are absorbed, and who can change key parameters. Fairness requires not only transparency but contestability: participants need some way to challenge exclusions, parameter changes, or suspicious operator behavior.

\textbf{Non-exploitative incentives.} The system should not depend on maximizing compulsive engagement. A framework meant to compensate attention contributors cannot justify itself by intensifying attention capture at any cost.


%% ========================================================================
\section{A Simple Model of Token Dynamics}
\label{sec:model}
%% ========================================================================

\subsection{Core Variables and Feedback Loop}
\label{sec:model-core}

ATEM's dynamics are driven by a reinforcing loop connecting four quantities: the number of active users $N$, effective attention $E$, revenue $R$, and the token price $P$. The loop operates as follows: more users generate more attention, more attention generates more revenue, more revenue funds larger buybacks, larger buybacks support higher token prices, and higher token prices attract more users.

We formalize these relationships with simplified functional forms. Real implementations would require empirical calibration; the value of the analysis lies in the qualitative relationships it reveals.

\textbf{Effective attention scales with network efficiency.} While human attention is biologically finite (bounded by an absolute time limit $h$ per user), the \emph{monetizable value} or \emph{matching efficiency} of that attention improves with scale. Richer content, better recommendations, and denser marketplaces mean that a unit of attention on a larger platform is more valuable than on a smaller one. We model aggregate effective attention $E(N)$ as:
\begin{equation}
E(N) = \mu(N) \cdot N \cdot h
\end{equation}
where $\mu(N)$ is a network-dependent efficiency multiplier. If we assume $\mu(N)$ scales as $N^{\alpha-1}$ (with $\alpha > 1$ capturing the strength of network effects), aggregate effective attention simplifies to:
\begin{equation}
E(N) = N^\alpha \cdot h
\label{eq:attention}
\end{equation}

\textbf{Revenue scales with effective attention.} Advertisers pay for access to efficient, highly targeted attention. Revenue is therefore:
\begin{equation}
R(N) = r \cdot E(N) = r \cdot h \cdot N^\alpha
\label{eq:revenue}
\end{equation}
where $r$ is revenue per unit of effective attention.

\textbf{Buybacks link revenue to token supply.} A fraction $\delta$ of revenue is allocated to purchasing tokens from the open market. Let $e$ denote the tokens emitted per user per period. At steady state, token supply is stable when emission equals absorption:
\begin{equation}
e \cdot N = \frac{\delta \cdot R(N)}{P}
\end{equation}

Solving for the steady-state token price:
\begin{equation}
P^* = \frac{\delta \cdot r \cdot h \cdot N^{\alpha-1}}{e}
\label{eq:price}
\end{equation}

This expression reveals a key insight: because the network multiplier makes aggregate attention increasingly efficient to monetize, the equilibrium token price grows as a power of the user base. Doubling the number of users more than doubles the token price, creating the positive feedback loop on which the system depends.

\textbf{User adoption responds to expected rewards.} Users participate when expected token rewards exceed their opportunity cost $c$:
\begin{equation}
e \cdot P^* > c
\end{equation}

\subsection{Critical Mass}
\label{sec:model-critical}

Substituting the steady-state price into the participation condition:
\begin{equation}
\delta \cdot r \cdot h \cdot N^{\alpha-1} > c
\end{equation}

This yields the critical mass threshold:
\begin{equation}
N^* = \left(\frac{c}{\delta \cdot r \cdot h}\right)^{1/(\alpha-1)}
\label{eq:critical}
\end{equation}

Below $N^*$, the system cannot generate sufficient revenue to support token prices that justify participation. Above $N^*$, the positive feedback loop becomes self-reinforcing. The critical mass depends inversely on the strength of network effects ($\alpha$), the revenue yield per attention unit ($r$), the biological attention bound ($h$), and the buyback fraction ($\delta$), and directly on the opportunity cost of participation ($c$).

This has a practical implication: ATEM implementations should prioritize features that strengthen network effects (increasing $\alpha$) and revenue per user (increasing $r$), not simply maximize user count.

\subsection{The Buyback Allocation Tradeoff}
\label{sec:model-buyback}

A platform must decide what fraction $\delta$ of revenue to allocate to token buybacks versus operational investment. The tradeoff has a clear structure: too little buyback and token price stagnates, weakening adoption; too much buyback and the platform starves itself of operational capital, degrading product quality and driving users away.

If we model platform quality as depending on operational spending, $Q = q \cdot (1 - \delta) \cdot R$, and user growth as proportional to both token attractiveness and platform quality, growth is maximized at an interior value of $\delta$. For a simplified growth function proportional to $\delta^\beta \cdot (1 - \delta)^\gamma$, the optimal buyback fraction is:
\begin{equation}
\delta^* = \frac{\beta}{\beta + \gamma}
\label{eq:optimal-delta}
\end{equation}
where $\beta$ captures the sensitivity of adoption to token price and $\gamma$ captures the sensitivity of retention to platform quality. When network effects are strong (high $\beta$), a larger share should flow to buybacks. The balance point will differ across implementations and growth stages.

\subsection{A Magnitude Estimate}
\label{sec:model-magnitude}

An important question is whether attention compensation can be economically meaningful or only symbolic. A rough estimate grounds the proposal in observable market parameters.

Global digital advertising revenue reached approximately \$750--800 billion in 2025, distributed across roughly 5.5 billion internet users \citep{emarketer2025}. This implies average per-user advertising revenue of approximately \$136--145 per year, though with enormous geographic variation: per-user ad revenue exceeds \$900 in the United States but falls below \$20 in many developing economies.

Not all of this revenue is available for contributor compensation. Consider an ATEM implementation operating in a specific vertical (e.g., e-commerce discovery or social media) with 100 million active users generating \$12 billion in annual advertising and commerce revenue, implying per-user revenue of \$120 per year, roughly in line with the global average. With a platform operating margin of 35\%, gross margin per user is approximately \$42. If 50\% of gross margin is directed to token buybacks, approximately \$21 per user per year flows to token absorption. (In terms of the formal model's $\delta$, which is defined as a fraction of total revenue, this corresponds to $\delta \approx 0.35 \times 0.5 = 0.175$. The magnitude estimate thus uses a more conservative effective buyback rate than the sensitivity analysis baseline of $\delta = 0.5$, which treats $\delta$ as a fraction of revenue directly.) If the platform also generates revenue from marketplace commissions, subscriptions, or premium services (as most large platforms do), the effective per-user revenue base would be larger, and buyback-funded compensation could reach \$30--50 per year. The range is illustrative rather than precise; actual figures would depend on the specific revenue mix and vertical.

These figures are modest in high-income economies. But three considerations affect their significance. First, the impact varies dramatically by region. In economies where median income is \$2,000--5,000 per year, \$30--50 in annual attention compensation represents 0.6--2.5\% of income. In lower-income regions where median per-capita income falls to \$500--800 per year, the same amount represents 3.8--10\% of annual income, a non-trivial supplement comparable to what micro-task and gig platforms deliver. Second, the estimate assumes a single vertical; participation across multiple ATEM implementations could multiply returns. Third, digital advertising markets continue to grow at 10--15\% annually, and the per-user revenue pool expands as more economic activity moves online.

The magnitude estimate confirms that attention compensation is meaningful as a supplement in specific contexts, particularly in lower-income economies and high-engagement verticals, rather than as a primary income replacement in the near term. This is consistent with the paper's framing of attention compensation as a complementary layer rather than a substitute for labor income.

\subsection{Sensitivity of Critical Mass}
\label{sec:model-sensitivity}

The critical mass formula~\eqref{eq:critical} depends on several parameters. A natural question is how sensitive $N^*$ is to plausible variation in the two most policy-relevant parameters: the network effect exponent $\alpha$ and the buyback fraction $\delta$. Table~\ref{tab:calibration} presents the baseline calibration, and Table~\ref{tab:sensitivity} shows the resulting $N^*$ across a grid of values.

\begin{table}[H]
\centering
\caption{Parameter calibration}
\label{tab:calibration}
\begin{tabularx}{\textwidth}{lccX}
\toprule
\textbf{Parameter} & \textbf{Symbol} & \textbf{Baseline} & \textbf{Interpretation} \\
\midrule
Network effect exponent & $\alpha$ & 1.15 & Weak to strong network effects (range: 1.05--1.30) \\
Buyback fraction & $\delta$ & 0.50 & Conservative to aggressive buyback (range: 0.20--0.70) \\
Revenue per effective attention & $r$ & calibrated & Market conditions (\$120/user/year at baseline) \\
Opportunity cost & $c$ & calibrated & Participation threshold ($N^* = 87$K at baseline) \\
\bottomrule
\end{tabularx}
\end{table}

\begin{table}[H]
\centering
\caption{Critical mass $N^*$ under varying $\alpha$ and $\delta$}
\label{tab:sensitivity}
\begin{tabular}{lrrrr}
\toprule
 & $\delta = 0.20$ & $\delta = 0.35$ & $\delta = 0.50$ & $\delta = 0.70$ \\
\midrule
$\alpha = 1.05$ & $> 10^{22}$ & $> 10^{17}$ & $> 10^{14}$ & $> 10^{11}$ \\
$\alpha = 1.10$ & $> 10^{11}$ & 908M & 25.7M & 887K \\
$\alpha = 1.15$ & 39.1M & 938K & \textbf{87K} & 9K \\
$\alpha = 1.30$ & 6K & 968 & 295 & 96 \\
\bottomrule
\end{tabular}
\end{table}

The baseline case ($\alpha = 1.15$, $\delta = 0.50$) yields $N^* = 87{,}000$, consistent with the calibration in Section~\ref{sec:model-critical}. Three observations emerge from the sensitivity analysis.

First, $N^*$ is extremely sensitive to $\alpha$. Moving from $\alpha = 1.15$ to $\alpha = 1.10$ (a modest reduction in network effects) increases critical mass by nearly 300x, from 87K to 25.7 million. Moving to $\alpha = 1.05$ pushes $N^*$ beyond $10^{11}$ even in the most favorable case, effectively making the system infeasible. This confirms that the viability of attention-token systems depends critically on whether attention aggregation exhibits meaningful network effects, not merely on user count or revenue.

Second, $N^*$ is moderately sensitive to $\delta$. At the baseline $\alpha = 1.15$, increasing $\delta$ from 0.20 to 0.70 reduces critical mass by a factor of roughly 4,000 (from 39.1 million to 9K). This is substantial but far less dramatic than the effect of $\alpha$, suggesting that the buyback fraction is a tunable policy lever rather than a make-or-break parameter.

Third, when network effects are strong ($\alpha = 1.30$), critical mass collapses to hundreds or low thousands of users, and the system becomes viable even with modest buyback allocations. This suggests that ATEM implementations in domains with inherently strong network effects (e.g., social platforms, marketplaces with strong cross-side effects) face fundamentally different viability challenges than those in domains with weak network effects.

These results should be interpreted with caution. The parameters are not independent in practice: $\alpha$ may itself vary with $N$, and $\delta$ is constrained by operating costs. The analysis assumes a single ATEM implementation in isolation (see Section~\ref{sec:model-limits} and Section~\ref{sec:limits-competition} for further discussion of multi-system dynamics). Nevertheless, the sensitivity pattern reinforces the paper's central practical recommendation: implementations should prioritize features that strengthen network effects over strategies that simply maximize user acquisition or buyback spending.

\subsection{Limitations of the Model}
\label{sec:model-limits}

This model deliberately trades precision for clarity. The functional forms are stylized; real attention markets exhibit heterogeneous user behavior, diminishing returns at very large scale, and competitive dynamics across platforms. The model analyzes one ATEM implementation in isolation; multiple competing systems would create different dynamics requiring game-theoretic analysis. The magnitude estimate relies on aggregate figures and cannot predict outcomes for specific implementations. The model is best understood as identifying what must be true for attention-based economies to work, rather than predicting exactly how they will unfold.


%% ========================================================================
\section{Limits, Risks, and Governance Problems}
\label{sec:limits}
%% ========================================================================

\subsection{External Demand and Macroeconomic Limits}
\label{sec:limits-demand}

An attention-compensation system cannot sustain itself without buyers. Advertisers, merchants, platforms, or other monetizers must have a reason to pay for verified human attention. A common objection is the macroeconomic paradox of a post-labor consumer base: if AI and automation severely reduce labor income, why would advertisers pay for attention from consumers with diminished purchasing power?

The answer lies in relative scarcity. As AI and automation expand, the ability to produce goods grows exponentially while production costs collapse toward zero \citep{rifkin2014}. Concurrently, human attention remains biologically fixed. This stark mismatch (an exponentially growing supply of zero-marginal-cost or deeply cheapened goods chasing a strictly finite amount of human attention) means that absolute commercial competition intensifies. Even as consumer income drops, the sheer volume of competing products ensures that the relative demand for human attention, and thus advertising spend, remains robust or even increases. Attention becomes the ultimate bottleneck of the economy.

A stronger version of the objection notes that advertisers do not pay for attention in the abstract; they pay for attention that converts to purchases. If consumer purchasing power declines sharply, advertising budgets should follow regardless of attention scarcity. This paper does not claim the paradox is fully resolved. Rather, it notes two partial responses. First, zero-marginal-cost goods (digital services, AI-generated content, software) can be ``purchased'' with attention alone, sustaining a commercial loop that does not require traditional purchasing power. Second, in a mixed economy where labor income declines gradually rather than collapses suddenly, attention compensation serves as a complementary income stream that partially sustains the purchasing power on which advertising markets depend. A full macroeconomic general equilibrium model of this dynamic is beyond the scope of this paper, but it remains one of the most important directions for future research.

\subsection{Surveillance, Proof-of-Humanity, and the AI Simulation Problem}
\label{sec:limits-surveillance}

The more valuable participation rewards become, the stronger the incentive to automate or fake them. Historically, platforms used behavioral proxies (browsing patterns, mouse movements, time-on-page) to filter bots. However, agentic AI introduces an existential threat to this approach: large language models and autonomous AI agents can now simulate ``human'' engagement behavior, including views, clicks, time-on-page, and natural-language interactions, rendering behavioral proxies obsolete and highly exploitable. Paradoxically, this makes verified human attention more rather than less valuable: it is precisely because AI can now simulate behavioral proxies that authentic, proof-of-humanity-verified attention becomes a scarcer and more premium resource for advertisers.

This forces a stark institutional dilemma. To stop advanced AI from farming an attention-token system, verification must inevitably lean toward stronger Proof-of-Humanity mechanisms \citep{borge2017}, such as invasive biometrics (e.g., Worldcoin's iris scanning) or highly centralized credential systems. But a framework built to fairly compensate attention cannot justify itself by mandating dystopian surveillance without clear limits and recourse. The tension between robust sybil resistance in an AI-dominated era and the ethical imperative of minimal surveillance is inherent rather than solvable. Implementations must make explicit where they sit on that spectrum and subject the choice to strict governance.

\subsection{Manipulation and Behavioral Dependency}
\label{sec:limits-manipulation}

Attention markets already have strong incentives to maximize time-on-platform, emotional arousal, and compulsive return behavior. A compensation layer could intensify those incentives if poorly governed. Designers therefore need explicit restraints on how rewards interact with platform optimization. The system should compensate attention that users would plausibly supply without addictive design, not reward users for being manipulated more effectively.

\subsection{Volatility, Liquidity, and Regulatory Risk}
\label{sec:limits-volatility}

Tokenized systems can make rewards visible and programmable, but they also expose participants to price volatility and uncertain regulation. This is especially problematic if the people targeted for inclusion are also the least able to absorb financial risk. Fairness requires confronting this directly. Possible mitigations include stablecoin-denominated rewards, minimum holding periods, and fiat conversion guarantees, but each introduces its own tradeoffs.

\subsection{Governance and Trust}
\label{sec:limits-governance}

Even if on-chain token movements are transparent, most monetization processes remain partly off-chain. Revenue attribution, merchant contracts, advertising quality, and operator decisions are difficult to verify in real time. No serious design can promise total trustlessness. It must specify what remains trust-based, what can be audited, and how participants can contest abuse. ATEM implementations may adopt various governance approaches, from algorithmic parameter adjustment to democratic participant voting to hybrid models. What matters is that governance be legitimate, transparent, and contestable.

\subsection{Lessons from Failed Token Systems}
\label{sec:limits-failures}

The recent history of token economies provides cautionary examples. Algorithmic stablecoin systems such as Terra/Luna collapsed when their stabilization mechanism, which relied on arbitrage between two linked tokens rather than on external revenue, entered a reflexive death spiral under selling pressure. Reward-token systems such as Steemit faced governance capture and inflationary distribution that decoupled rewards from any measure of real demand.

ATEM's architecture differs in a structural respect: buyback pressure derives from external revenue (advertising, commerce, subscriptions) rather than from internal token arbitrage or speculative inflows. This does not make ATEM immune to failure. If external revenue declines, the buyback mechanism weakens proportionally. But the failure mode is contraction rather than reflexive collapse, because the system does not depend on internal token dynamics to maintain its own stability. The discipline comes from outside the token system, through the willingness of advertisers and merchants to pay for verified human attention.

\subsection{Revenue Circularity and Competition}
\label{sec:limits-competition}

The model assumes an external real economy generating advertising and commerce demand. It also analyzes one ATEM implementation in isolation. Multiple competing attention-token systems would fragment user bases and weaken network effects for each, potentially pushing all systems below their critical mass thresholds. Whether attention-compensation systems will tend toward consolidation (like social media platforms) or fragmentation (like cryptocurrency markets) remains an open empirical question.


%% ========================================================================
\section{Conclusion}
\label{sec:conclusion}
%% ========================================================================

This paper argues that human attention should be recognized as a form of economic participation in domains where it is directly monetizable, and that such recognition becomes more pressing as AI and automation make labor-centered income less secure.

The paper makes three contributions. First, the Attention Triad distinguishes contribution, aggregation, and monetization as separate economic functions in attention markets, making the contribution gap visible as a structural feature: contributors supply the scarce input, yet usually remain outside the direct distributive loop. The Triad also clarifies how attention-based compensation relates to and differs from the data-as-labor framework proposed by \citet{posner2018} and \citet{arrieta2018}.

Second, the attention-centered economy framework identifies a structural parallel between labor-centered and attention-centered economic arrangements and articulates the bridging insight that the same forces weakening labor-based inclusion simultaneously strengthen the economic centrality of human attention. This connection, previously undrawn in the literature, provides a principled basis for taking attention compensation seriously as a distributive mechanism rather than dismissing it as a platform marketing device.

Third, ATEM specifies one governance-conscious institutional architecture for returning part of the value generated by attention markets to verified human participants, along with a simple formal model that identifies the structural conditions for viability: a critical mass threshold inversely related to network effects, a buyback allocation tradeoff that determines system growth, and a magnitude estimate showing that attention compensation is meaningful as a complement in specific contexts rather than as a primary income replacement. The sensitivity analysis (Section~\ref{sec:model-sensitivity}) further demonstrates that viability depends critically on network effects rather than on any single policy lever, reinforcing the importance of institutional design over parameter tuning.

These contributions carry implications for policy. Regulators considering platform compensation mandates can use the Attention Triad to distinguish between genuine attention compensation (structurally linked to revenue and governed by contestable rules) and marketing relabeled as user rewards. Policymakers designing digital inclusion or UBI programs may find attention compensation a complementary mechanism that operates within markets rather than relying solely on fiscal redistribution. Tax and securities regulators face an emerging question about how attention-compensation tokens should be classified: as income, as securities, as utility tokens, or as a category that existing frameworks do not yet address.

That proposal remains conditional. A fair attention-compensation system would need real-demand linkage, anti-sybil protection, privacy restraint, contestable governance, non-exploitative incentives, and a more credible account of what forms of attention buyers actually value. Revenue-backed token sinks can help discipline the system, but they do not guarantee sustainability. A tokenized implementation may help operationalize ATEM, but the underlying normative claim is broader: contributors should have a governed and intelligible claim on value where their attention is already being monetized.

These limits do not make the problem less important. If AI and automation continue to weaken labor's role as the ordinary person's primary channel of economic participation, then the undercompensation of attention contributors will become harder to treat as a marginal issue. The next step is not to overstate what attention can bear. It is to build a more rigorous, empirically grounded literature on when attention compensation is fair, legitimate, and institutionally defensible.


%% ========================================================================
%% References
%% ========================================================================

\bibliographystyle{apalike}
\bibliography{references}

\end{document}
